\documentclass[11pt]{article}
\usepackage[margin=1in]{geometry}
\usepackage{hyperref}
\usepackage{enumitem}
\usepackage{booktabs}
\usepackage{xcolor}
\usepackage{amsmath}
\title{ProbOS --- Mission, Vision, and Industry Roadmap}
\author{Nisong Monyimba \\ Reality Computing Corporation}
\date{\today}

\begin{document}
\maketitle

\section*{A note on how to read this document}
This document states ProbOS's mission and long-term vision, then
sequences which industries it will serve and in what order, honestly
distinguishing three categories throughout:
\begin{itemize}[leftmargin=*]
  \item \textbf{Validated} -- backed by actual completed work
    (a literature-validated model, a passing test, a measured
    number), not just plausibility.
  \item \textbf{Structurally ready} -- the underlying kernel already
    supports this with little or no new engineering, but the
    domain-specific model itself has not been built or validated yet.
  \item \textbf{Aspirational} -- a real, considered target, but with
    no work done toward it yet and no practitioner conversation
    confirming the need.
  \end{itemize}
No claim in this document is marked Validated unless it corresponds
to something that actually exists and has been tested, per
\texttt{docs/standards/quality\_standards.md}.

\section*{Mission}
ProbOS treats uncertainty as a first-class data type in software that
models safety-critical physical systems. Where deterministic
simulation uses one parameter value where there should be a
distribution, ProbOS propagates the full range of physically
plausible values through a model and reports the tail risk that a
single-point simulation cannot see.

\section*{Vision}
A general-purpose, open, auditable probabilistic execution runtime
that becomes the default tool for engineers who must not only
simulate a safety-critical system, but \emph{prove to a regulator}
that they understand its failure modes -- with a causal record from
every dangerous predicted outcome back to the specific parameters and
assumptions that produced it.

\section*{The core thesis: why safety-critical UQ, specifically}
ProbOS does not attempt to be a general replacement for existing
probabilistic programming frameworks (Stan, PyMC, NumPyro, Pyro,
Turing.jl), which are mature, widely used, and well-suited to
statistical inference and machine learning workflows. ProbOS's
distinct positioning rests on four criteria, present together in a
specific, underserved niche:
\begin{enumerate}[leftmargin=*]
  \item A formal regulator requires documented failure-mode analysis.
  \item The system is modelable as an ODE/SDE with genuine,
    quantifiable parameter uncertainty (material properties,
    manufacturing variance, patient variability).
  \item Failure consequences are rare but catastrophic -- exactly the
    tail risk a deterministic point-estimate simulation is structurally
    blind to.
  \item Existing deterministic tools in that industry are documented
    to miss these tail risks.
\end{enumerate}
This is a narrower, more specific claim than ``uncertainty
quantification is useful'' -- it identifies exactly which industries
have this combination of properties, and treats industries lacking
one or more of them as a poor fit, not a target to chase for breadth.

\section*{Validated today}

\subsection*{Battery safety (EV, grid storage) -- Validated}
\texttt{BatteryModel2Cell}, an 8-state Arrhenius thermal-runaway ODE
with 15 uncertain parameters, validated against Kim et al. (2007)
accelerating rate calorimetry data (Month 1 Week 2). Sobol sensitivity
analysis (Month 1 Week 3) identifies \texttt{Ea\_SEI} as the dominant
driver of thermal-runaway variance ($S_1 = 0.457$), and Monte Carlo
propagation shows the P95 worst-case cell decomposes 11.6x faster than
the P50 median -- a tail risk invisible to a deterministic simulation
using literature-mean parameters. The provenance tracker (Month 1
Week 3) traces any dangerous predicted output back to the specific
parameter draws that produced it, directly serving the ``documented
failure-mode analysis'' requirement regulators impose.

\section*{Structurally ready -- near-zero additional engineering}

\subsection*{Pharmaceutical stability / shelf-life (FDA CMC, ICH Q1A) -- Validated}
\texttt{PharmaStabilityModel} (Month 3 Week 11) uses Avrami kinetics
(a related but distinct reaction-kinetics law from
\texttt{BatteryModel2Cell}'s first-order depletion, since the
real, best-validated dataset found required it), validated against
Gonzalez-Gonzalez et al. (\emph{Molecules} 2023, 28(23), 7925), a
real, peer-reviewed, ICH-referenced chlorhexidine stability study.
The activation energy (Ea = 18.52 kcal/mol) is independently cited
and cross-checked against an independent 1987 PhD thesis range
(16.5--22.9 kcal/mol). The model reproduces the paper's 3 real
reported 365-day degradation percentages (5C/25C/30C) to within
1.18 percentage points.

\textbf{Honest methodological caveat:} unlike Kim (2007)'s true
forward replication for the battery model, this validation fits
the pre-exponential factor A against the paper's own reported
summary results, with the Avrami exponent n fixed at 1.0 as an
explicit simplifying assumption (confirmed necessary: an
unconstrained fit of both A and n against only 3 single-timepoint
data points is fundamentally underdetermined, and converged to a
physically absurd result when attempted directly). This is a real
but weaker validation standard than Kim (2007) and should not be
presented as equivalent rigor.

Sobol sensitivity analysis (using the paper's own reported Ea
standard error, $\pm$2.61 kcal/mol, as a genuine literature-backed
uncertainty) identifies Ea as overwhelmingly dominant
($S_1 = 0.993$), and Monte Carlo propagation reveals a real tail
risk invisible to deterministic simulation: the worst-case 5\% of
outcomes show complete potency loss within one year, versus 82.6\%
potency remaining at the deterministic point estimate.

\subsection*{Automotive functional safety / EV battery management (ISO 26262) -- Structurally ready}
\texttt{BatteryModel2Cell} (already Validated) is positioned, not
re-engineered, for this market. Month 3 Week 12 produced accurate
positioning materials
(\texttt{docs/automotive\_ev\_positioning.md},
\texttt{docs/automotive\_ev\_outreach\_email.md}) grounded in real
ISO 26262 research (HARA, Part 3 Clause 7). An honest, load-bearing
distinction is maintained throughout: ProbOS does NOT perform HARA,
does NOT assign ASIL levels, and has NOT been used in an actual
ISO 26262 submission. It provides quantitative uncertainty evidence
(Sobol-identified dominant risk drivers, Monte Carlo tail-risk
quantification, provenance-traced causal explanations) that could
genuinely support a HARA process's Severity/Exposure judgments --
materials exist; the actual outreach conversation has not yet
happened.

\section*{Aspirational -- real targets, no work done yet}

\subsection*{Medical device (implantable battery safety, FDA 510k) -- Aspirational}
Named in \texttt{docs/sbir\_phase1\_onepager.md} as an SBIR Phase I
target (FDA Q-Sub filing for implantable battery MRI safety). No
domain-specific model has been built or validated for this industry
yet; the provenance tracker's audit-trail capability is directly
relevant once a specific device model exists.

\subsection*{Nuclear safety analysis (NRC) -- Aspirational}
Named as a target market throughout project documentation since Month
1. No nuclear-specific model exists. NRC's uncertainty quantification
requirements match the core thesis structurally, but require a new
physical model (reactor thermal/neutronics or similar) not yet
attempted.

\subsection*{Aerospace / aviation (FAA, DO-178C / DO-254) -- Aspirational}
Structurally the closest sibling to the current battery/medical/
nuclear framing: same regulatory pattern (documented, auditable
failure analysis), same tail-risk stakes. Would require new
domain models (structural fatigue, engine thermal, or flight-control
parameter uncertainty) -- none built yet.

\subsection*{Grid / power systems reliability (NERC) -- Aspirational, higher effort}
A real and growing need given renewable-generation variability
introducing genuine parameter uncertainty into grid stability models,
and NERC's post-incident investigation requirements resembling the
provenance/audit-trail pitch already made for medical devices. Rated
higher-effort than the above because it requires a new class of
stochastic model (demand/generation processes), not an ODE-based
physical model like the others.

\subsection*{Explicitly out of current scope}
Quantitative finance and robotics are named in earlier project
documentation as market categories, but honest reassessment
(Month 3) concluded neither is currently a good fit: the finance
example model (\texttt{OptionPricerModel}, Black-Scholes GBM) is a
toy validation case rather than something serving a real quant
research audience already well-served by mature tools, and no
robotics-specific model exists at all. These are removed from active
targeting, not merely deprioritized.

\section*{Month-by-month roadmap}

\subsection*{Completed (Months 1--3, Weeks 1--10)}
\begin{itemize}[leftmargin=*]
  \item \textbf{Month 1 (Weeks 1--4):} Kernel foundation --
    \texttt{Distribution}/\texttt{Model} ABCs, \texttt{BatteryModel2Cell}
    validated against Kim (2007), \texttt{MonteCarloEngine},
    \texttt{SobolSensitivity}, \texttt{ProvenanceTracker}, C++/OpenMP
    kernel (7x measured speedup), PDSL v0.1, three cross-discipline
    forward-simulation examples.
  \item \textbf{Month 2 (Weeks 5--8):} \texttt{ParticleFilter}
    (validated against an exact closed-form Kalman filter), pybind11
    bindings, FastAPI service layer, cross-discipline sequential
    filtering (ED queue, clinical trial), plus a comprehensive
    mid-month hardening pass closing several real, previously
    undiscovered gaps (mypy/ruff scope, a license mismatch, vestigial
    code).
  \item \textbf{Month 3 (Weeks 9--10, in progress):} REST API model
    registry extended to all four models; GPU kernel path built,
    honestly benchmarked, and root-cause-optimized via kernel fusion
    -- C++/OpenMP remains the fastest engine for this workload on
    this hardware, reported truthfully rather than oversold.
\end{itemize}

\subsection*{Near-term (rest of Month 3 / Month 4) -- Structurally ready extensions}
\begin{itemize}[leftmargin=*]
  \item Pharmaceutical stability model, reusing
    \texttt{BatteryModel2Cell}'s Arrhenius architecture with new
    priors, validated against a real stability-study dataset.
  \item Automotive/EV safety positioning and outreach, using the
    existing validated battery model under an ISO 26262 framing.
  \item PDSL v0.2 (control flow), continuing the priority order
    already established in Month 3's original plan.
\end{itemize}

\subsection*{Mid-term (Months 5--6) -- Validation, not more building}
\begin{itemize}[leftmargin=*]
  \item At least one real conversation with a practicing battery,
    pharmaceutical, or automotive safety engineer, confirming
    (or correcting) the assumed pain point and whether ProbOS's
    current capabilities actually address it.
  \item Medical device domain model, informed by whatever is learned
    from the above conversations rather than built speculatively in
    advance of them.
\end{itemize}

\subsection*{Longer-term (Month 7+) -- Contingent on validated demand}
\begin{itemize}[leftmargin=*]
  \item Nuclear and aerospace domain models, pursued only if the
    mid-term validation step confirms genuine practitioner interest
    in this pattern -- not built speculatively for market-size
    breadth.
  \item Grid/power systems reliability, the highest-effort item on
    this list, deferred until the lower-effort extensions have
    actually been validated.
\end{itemize}

\section*{Honest open questions}
This document deliberately does not claim ProbOS is close to
displacing established tools, or that its market thesis is confirmed.
Two open questions remain genuinely unresolved as of this writing:
\begin{enumerate}[leftmargin=*]
  \item \textbf{Competitive differentiation.} What, specifically, does
    ProbOS do that Stan, PyMC, or NumPyro do not, for an engineer who
    already has access to those tools? The provenance/audit-trail
    angle is the strongest current answer, but has not been tested
    against a real user's actual workflow.
  \item \textbf{Practitioner validation.} Every domain claim in this
    document, including the Validated battery-safety work, has been
    checked against literature data, not against a live conversation
    with someone who would actually adopt the tool. This is the
    single highest-priority open item ahead of further domain
    expansion.
\end{enumerate}

\end{document}
